\section{Function guide}
\label{sec:api-reference}

The functions below cover most everyday calculations. API means
\emph{application programming interface}: the public classes and
functions you call from Python. The complete reference, including
less frequently used utilities and every exact signature, is in
\path{docs/api.md}.

\subsection{The functions you will use most}

\begin{tabular}{@{}p{0.39\linewidth}p{0.56\linewidth}@{}}
\toprule
\textbf{Function or class} & \textbf{When to use it} \\
\midrule
\texttt{SpectroscopySystem} & Define the Hamiltonian, dipole, dissipation
and optional initial state. \\
\texttt{GaussianPulse} & Set a Gaussian pulse's carrier frequency, temporal
width, amplitude and phase. \\
\texttt{system.stationary\_state()} & Obtain a stationary initial density
matrix; pass it as \texttt{rho0} in the calculation. \\
\texttt{*\_spectrum} & Calculate the frequency array you will normally
plot: a 1D linear spectrum or a 2D third-order spectrum. \\
\texttt{plot\_linear\_spectrum} & Plot one component of a 1D linear spectrum;
pass an existing axis to overlay curves. \\
\texttt{plot\_spectrum} & Plot one component of a 2D spectrum on its
Cartesian frequency grid. \\
\texttt{finite\_pulse\_points} & Calculate selected frequency pairs
without building a full 2D grid. \\
\bottomrule
\end{tabular}

\subsection{What do signal, spectrum, and response mean?}

Choose a model prefix, such as \texttt{finite\_pulse} or
\texttt{impulsive\_rwa}; its suffix tells you how to supply inputs
and what is returned:
\begin{description}
\item[\texttt{\_spectrum}.] Use a system object and frequency arrays.
The result is a complex array: \texttt{(N,)} for linear spectroscopy,
or \texttt{(N3, N1)} for third order.
\item[\texttt{\_signal}.] Use a system object and a single frequency
coordinate. The result is one complex number. Use this for one peak or
one selected pair rather than a complete plot.
\item[\texttt{\_response}.] Supply the Liouvillian, interaction and
detection matrices directly. The result is one complex number.
These functions are useful when working directly with superoperators;
most users can use \texttt{\_signal} or \texttt{\_spectrum} instead.
\item[\texttt{\_time\_signal}.] Supply delay-time arrays instead of
frequencies. The result is a complex time-domain array. All such functions
in this version use impulsive interactions.
\end{description}

These are interface names: \texttt{signal} does not add a detector model.
For example, the finite-pulse NR response can be evaluated at one pair
or across a full Cartesian grid:
\par\noindent
\begin{minipage}{\linewidth}
\begin{lstlisting}[style=python]
value = pulsus.finite_pulse_signal(
    system, pulse, pulse, pulse, 1.24, 1.24,
    T=20.0, eta=0.02, pathway="NR", rho0=rho_ss,
)  # one complex number
S = pulsus.finite_pulse_spectrum(
    system, pulse, pulse, pulse, omega1, omega3,
    T=20.0, eta=0.02, pathway="NR", rho0=rho_ss,
)  # shape (len(omega3), len(omega1))
\end{lstlisting}
\end{minipage}\par

\clearpage
\subsection{Common calling patterns}

The three third-order Gaussian-pulse models use the same arguments:
\par\noindent
\begin{minipage}{\linewidth}
\begin{lstlisting}[style=python]
finite_pulse_spectrum(
    system, pulse1, pulse2, pulse3, omega1, omega3, T, eta,
    pathway="NR", rho0=None, signature=None)
finite_pulse_rwa_spectrum(
    system, pulse1, pulse2, pulse3, omega1, omega3, T, eta,
    pathway="NR", rho0=None, signature=None)
short_pulse_rwa_spectrum(
    system, pulse1, pulse2, pulse3, omega1, omega3, T, eta,
    pathway="NR", rho0=None, signature=None)
\end{lstlisting}
\end{minipage}\par
The impulsive models take no pulse object:
\par\noindent
\begin{minipage}{\linewidth}
\begin{lstlisting}[style=python]
impulsive_spectrum(
    system, omega1, omega3, T, eta, rho0=None)
impulsive_rwa_spectrum(
    system, omega1, omega3, T, eta,
    pathway="NR", rho0=None, signature=None)
\end{lstlisting}
\end{minipage}\par
Replacing \texttt{\_spectrum} with \texttt{\_signal} in any of these
five names keeps the arguments but requires scalar frequencies.
\texttt{finite\_pulse\_points} uses the finite-pulse signature with
paired arrays of equal length and returns an \texttt{(N,)} array.

For linear spectroscopy, the four spectrum functions are
\par\noindent
\begin{minipage}{\linewidth}
\begin{lstlisting}[style=python]
finite_pulse_linear_spectrum(
    system, pulse, omega, eta, sign=1, rho0=None)
short_pulse_linear_rwa_spectrum(
    system, pulse, omega, eta, sign=1, rho0=None)
impulsive_linear_spectrum(system, omega, eta, rho0=None)
impulsive_linear_rwa_spectrum(
    system, omega, eta, sign=1, rho0=None)
\end{lstlisting}
\end{minipage}\par
Here \texttt{omega} is one 1D axis. Replacing \texttt{\_spectrum} with
\texttt{\_signal} takes one frequency and returns one complex number.
There is no separate finite-pulse linear RWA family in this version.

\subsection{Plot, customize and optionally save}
\label{sec:plot-export}

Both plotting helpers return Matplotlib objects. They create a figure if
\texttt{ax} is omitted, and draw on that axis if it is supplied:
\par\noindent
\begin{minipage}{\linewidth}
\begin{lstlisting}[style=python]
linear_spectrum = spectra["Finite pulse, full"]
fig, ax, line = pulsus.plot_linear_spectrum(
    omega, linear_spectrum, component="imag", normalize=True,
)
plt.show()

fig, ax, contour = pulsus.plot_spectrum(
    omega1, omega3, S, component="imag", normalize=True,
)
plt.show()
\end{lstlisting}
\end{minipage}\par
Use \texttt{component="real"}, \texttt{"imag"}, or \texttt{"abs"}.
Normalization scales only the displayed component; it does not modify
the complex input array. The linear helper accepts ordinary line options,
such as \texttt{linestyle="--"}, \texttt{color="navy"} and a curve
\texttt{label}; add \texttt{ax.legend()} when comparing labeled curves.
Neither helper displays or saves automatically. Use \texttt{plt.show()}
when running a script.

Both helpers accept \texttt{savepath=None} and \texttt{dpi=300}.
Omitting \texttt{savepath} writes nothing. Supply a filename to export the
figure immediately; its extension selects the format:
\begin{lstlisting}[style=python]
fig, ax, line = pulsus.plot_linear_spectrum(
    omega, linear_spectrum, component="imag",
    savepath="linear_spectrum.png", dpi=300,
)
# Use savepath="linear_spectrum.pdf" for PDF output.
\end{lstlisting}
The same two optional arguments work with \texttt{plot\_spectrum}.
Use an existing destination folder. The resolution controls raster
output; PDF line work remains vector graphics.

For additional curves, legends or other customization, omit
\texttt{savepath} and export the completed figure afterward:
\begin{lstlisting}[style=python]
ax.set_title("My spectrum")
fig.savefig("custom_spectrum.pdf", dpi=300, bbox_inches="tight")
plt.show()
\end{lstlisting}
Saving a figure exports its visual appearance, not the complex numerical
spectrum. The input array remains available for analysis and separate data
storage.

Refer to Section~\ref{sec:time-domain} for time-domain
functions and to \path{docs/api.md} for lower-level response functions,
Liouville-space utilities, resolvents, and Gaussian dressing functions.
